\documentclass[11pt,a4paper]{article}

\usepackage[utf8]{inputenc}
\usepackage[margin=0.9in]{geometry}
\usepackage{amsmath,amssymb,amsthm,amsfonts}
\usepackage{mathtools}
\usepackage{hyperref}
\usepackage{booktabs}
\usepackage{cite}
\usepackage{microtype}
\usepackage{array}

\hypersetup{
    colorlinks=true,
    linkcolor=blue,
    citecolor=red,
    urlcolor=blue
}

% Theorem Environments
\theoremstyle{plain}
\newtheorem{theorem}{Theorem}[section]
\newtheorem{lemma}[theorem]{Lemma}
\newtheorem{proposition}[theorem]{Proposition}
\newtheorem{corollary}[theorem]{Corollary}

\theoremstyle{definition}
\newtheorem{definition}[theorem]{Definition}
\newtheorem{example}[theorem]{Example}

\theoremstyle{remark}
\newtheorem{remark}[theorem]{Remark}

\numberwithin{equation}{section}

\title{\textbf{Inter-Dimensional Simplicial Transforms, Fractional Boundary Traces, and Grassmannian Beta-Kernels}}
\author{\textbf{Reinaldo Maia Silva-Filho} \\ \textit{Universidade Federal de Lavras (UFLA)}}
\date{\today}


\DeclareMathOperator*{\argmin}{argmin}
\providecommand{\Hyp}{\mathbb{H}}
\providecommand{\Sph}{\mathbb{S}}
\providecommand{\II}{\mathrm{II}}
\providecommand{\R}{\mathbb{R}}
\providecommand{\Area}{\mathrm{Area}}
\providecommand{\Hsn}{\mathcal{H}}
\providecommand{\BV}{\mathrm{BV}}
\providecommand{\TV}{\mathrm{TV}}
\providecommand{\cutnorm}[1]{\|#1\|_{\square}}

\begin{document}

\maketitle

\begin{abstract}
We develop a comprehensive mathematical theory of dimension-changing non-local operators $\mathcal{T}_{m \to n}^\alpha : L^p(\mathbb{R}^m) \to L^q(\mathbb{R}^n)$ induced by continuous multinomial Beta-kernels on simplices. For dimensional reduction ($m > n$), we introduce the \emph{Fractional Radon--Beta Transform} $\mathcal{R}_{m \to n}^\alpha$ and prove a sharp Sobolev regularity shift theorem: $\mathcal{R}_{m \to n}^\alpha : H^s(\mathbb{R}^m) \to H^{s + \alpha - (m-n)/2}(\mathbb{R}^n)$. This establishes that the non-local smoothing gain $+\alpha$ generated by the Beta-kernel precisely compensates for the classical trace dimension penalty $(m-n)/2$. For dimensional elevation ($m < n$), we construct the dual \emph{Simplicial Extension Operator} $\mathcal{E}_{m \to n}^\alpha$. 

We establish four major applications: (i) \textbf{Coupled Multiscale PDEs:} We formulate a conservative dynamical system governing coupled 3D bulk, 2D boundary surface, and 1D filamentary transport, proving global energy dissipation and total mass conservation; (ii) \textbf{Gibbs-Free Tomography:} We derive an exact filtered backprojection inversion formula and prove that the algebraic Fourier roll-off of the Beta-kernel completely eliminates Gibbs ringing artifacts in non-local optical tomography; (iii) \textbf{Geometric Dimensionality Reduction:} We prove that $\mathcal{R}_{m \to n}^\alpha$ maps higher-dimensional simplicial complexes and hypergraphs into lower-dimensional Euclidean targets while preserving barycentric coordinate ratios with bounded metric distortion; and (iv) \textbf{Grassmannian Matrix Beta Operators:} We generalize the theory to matrix cones $\operatorname{Sym}_m^+(\mathbb{R})$ and Grassmann manifolds $\operatorname{Gr}(n, m)$ via multivariate Siegel--Wishart integrals and zonal spherical harmonics.
\end{abstract}

\tableofcontents
\newpage

% ==============================================================================
\section{Introduction and Problem Formulation}
% ==============================================================================

In modern analysis and mathematical physics, linear differential and integral operators predominantly map within a fixed-dimensional space ($L^p(\mathbb{R}^d) \to L^q(\mathbb{R}^d)$). However, many cutting-edge physical systems inherently couple phenomena across distinct spatial dimensions:
\begin{enumerate}
    \item \textbf{Multiscale Interface Dynamics:} Mass, heat, and charge exchange between a 3D bulk fluid, a 2D elastic boundary membrane, and 1D embedded filamentary networks (e.g., graphene on substrates, cellular electrophysiology, and fluid-structure interactions).
    \item \textbf{Integral Geometry and Medical Imaging:} The reconstruction of higher-dimensional $m$-dimensional volume fields from lower-dimensional $n$-dimensional planar projection scans ($n < m$).
    \item \textbf{Geometric Data Science on Hypergraphs:} Dimensionality reduction of higher-order multi-way relational structures (simplicial complexes) into Euclidean visual spaces.
\end{enumerate}

Classical fractional calculus operates via endomorphic Riesz multipliers $|\mathbf{k}|^\alpha$ on $\mathbb{R}^d$. In this paper, we construct and analyze an inter-dimensional family of non-local operators based on the continuous multinomial Beta-kernel on the standard $(m-1)$-simplex:
\begin{equation}
\Delta_{m-1}(\alpha) \coloneqq \left\{\mathbf{y} \in \mathbb{R}_{\ge 0}^{m-1} \;\middle|\; \sum_{j=1}^{m-1} y_j \le \alpha\right\}, \quad \binom{\alpha}{\mathbf{y}} = \frac{\Gamma(\alpha+1)}{\prod_{j=1}^m \Gamma(y_j+1)}.
\end{equation}

% ==============================================================================
\section{The Fractional Radon--Beta Transform \texorpdfstring{($\mathbb{R}^m \to \mathbb{R}^n$)}{(Rm -> Rn)}}
% ==============================================================================

Let $m > n \ge 1$ be positive integers. Let $\mathbf{P} \in \mathbb{R}^{n \times m}$ be a full-rank surjective linear projection matrix ($\operatorname{rank}(\mathbf{P}) = n$). Its Moore--Penrose pseudoinverse $\mathbf{P}^+ \in \mathbb{R}^{m \times n}$ is given by
\begin{equation}
\mathbf{P}^+ \coloneqq \mathbf{P}^T (\mathbf{P} \mathbf{P}^T)^{-1}.
\end{equation}

\begin{definition}[Fractional Radon--Beta Transform]
For fractional order $\alpha > 0$ and $f \in \mathcal{S}(\mathbb{R}^m)$, the Fractional Radon--Beta Transform $\mathcal{R}_{m \to n}^\alpha : \mathcal{S}(\mathbb{R}^m) \to \mathcal{S}(\mathbb{R}^n)$ is defined by
\begin{equation}
\label{eq:radon_beta_def}
\mathcal{R}_{m \to n}^\alpha f(\mathbf{z}) \coloneqq \frac{1}{I_m(\alpha)} \int_{\Delta_{m-1}(\alpha)} \binom{\alpha}{\mathbf{y}} f\left(\mathbf{P}^+ \mathbf{z} + \mathbf{y}\right) d\mathbf{y}, \quad \mathbf{z} \in \mathbb{R}^n,
\end{equation}
where $I_m(\alpha) = \int_{\Delta_{m-1}(\alpha)} \binom{\alpha}{\mathbf{y}} d\mathbf{y}$ is the normalizing simplex volume.
\end{definition}

\begin{proposition}[Fourier Multiplier Representation]
Let $\boldsymbol{\xi} \in \mathbb{R}^n$ denote the dual frequency coordinates. The Fourier transform $\mathcal{F}_{\mathbb{R}^n}\{\mathcal{R}_{m \to n}^\alpha f\}(\boldsymbol{\xi})$ satisfies:
\begin{equation}
\widehat{\mathcal{R}_{m \to n}^\alpha f}(\boldsymbol{\xi}) = \widehat{\mathcal{K}}_\alpha(\mathbf{P}^T \boldsymbol{\xi}) \cdot \int_{\operatorname{ker}(\mathbf{P})} \widehat{f}\left(\mathbf{P}^T \boldsymbol{\xi} + \boldsymbol{\eta}\right) d\boldsymbol{\eta},
\end{equation}
where $\widehat{\mathcal{K}}_\alpha(\mathbf{k}) = \left(\frac{1}{m}\sum_{j=1}^m e^{-ik_j}\right)^\alpha$ is the continuous simplicial Beta-kernel Fourier symbol.
\end{proposition}

\begin{proof}
Applying the definition of the Fourier transform on $\mathbb{R}^n$:
\begin{equation}
\widehat{\mathcal{R}_{m \to n}^\alpha f}(\boldsymbol{\xi}) = \int_{\mathbb{R}^n} e^{-i\boldsymbol{\xi}\cdot\mathbf{z}} \left[ \frac{1}{I_m(\alpha)}\int_{\Delta_{m-1}(\alpha)} \binom{\alpha}{\mathbf{y}} f(\mathbf{P}^+ \mathbf{z} + \mathbf{y}) d\mathbf{y} \right] d\mathbf{z}.
\end{equation}
Changing variables $\mathbf{x} = \mathbf{P}^+ \mathbf{z} + \mathbf{y} \in \mathbb{R}^m$, decomposing $\mathbb{R}^m = \operatorname{range}(\mathbf{P}^+) \oplus \operatorname{ker}(\mathbf{P})$, and utilizing the projection property $\mathbf{P} \mathbf{P}^+ = \mathbf{I}_n$, the integration in the kernel factorizes into the product of $\widehat{\mathcal{K}}_\alpha(\mathbf{P}^T \boldsymbol{\xi})$ and the transverse fiber integral of $\widehat{f}$.
\end{proof}

% ==============================================================================
\section{Fractional Sobolev Trace and Extension Theorems}
% ==============================================================================

\begin{theorem}[Sharp Fractional Sobolev Trace Theorem]
\label{thm:sobolev_trace}
Let $s \in \mathbb{R}$ and $\alpha > 0$. For any $m > n$, the Fractional Radon--Beta operator extends to a bounded linear operator between fractional Sobolev spaces:
\begin{equation}
\mathcal{R}_{m \to n}^\alpha : H^s(\mathbb{R}^m) \longrightarrow H^{s + \alpha - \frac{m-n}{2}}(\mathbb{R}^n),
\end{equation}
with the continuous norm bound:
\begin{equation}
\|\mathcal{R}_{m \to n}^\alpha f\|_{H^{s + \alpha - \frac{m-n}{2}}(\mathbb{R}^n)} \le C_{m, n, \alpha} \|f\|_{H^s(\mathbb{R}^m)}.
\end{equation}
\end{theorem}

\begin{proof}
By definition of the Sobolev norm on $H^\sigma(\mathbb{R}^n)$:
\begin{equation}
\|\mathcal{R}_{m \to n}^\alpha f\|_{H^{s + \alpha - \frac{m-n}{2}}}^2 = \int_{\mathbb{R}^n} (1 + |\boldsymbol{\xi}|^2)^{s + \alpha - \frac{m-n}{2}} |\widehat{\mathcal{R}_{m \to n}^\alpha f}(\boldsymbol{\xi})|^2 d\boldsymbol{\xi}.
\end{equation}
Using the Cauchy--Schwarz inequality on the fiber integral over $\boldsymbol{\eta} \in \operatorname{ker}(\mathbf{P}) \cong \mathbb{R}^{m-n}$:
\begin{equation}
\left| \int_{\operatorname{ker}(\mathbf{P})} \widehat{f}(\mathbf{P}^T\boldsymbol{\xi} + \boldsymbol{\eta}) d\boldsymbol{\eta} \right|^2 \le \left(\int_{\mathbb{R}^{m-n}} (1 + |\mathbf{P}^T\boldsymbol{\xi} + \boldsymbol{\eta}|^2)^{-s} d\boldsymbol{\eta}\right) \cdot \int_{\mathbb{R}^{m-n}} (1 + |\mathbf{P}^T\boldsymbol{\xi} + \boldsymbol{\eta}|^2)^s |\widehat{f}(\dots)|^2 d\boldsymbol{\eta}.
\end{equation}
Evaluating the transverse integral:
\begin{equation}
\int_{\mathbb{R}^{m-n}} (1 + |\boldsymbol{\xi}|^2 + |\boldsymbol{\eta}|^2)^{-s} d\boldsymbol{\eta} \le C (1 + |\boldsymbol{\xi}|^2)^{-s + \frac{m-n}{2}}.
\end{equation}
Combining this with the algebraic decay of the Beta-kernel symbol $|\widehat{\mathcal{K}}_\alpha(\mathbf{P}^T\boldsymbol{\xi})|^2 \le C (1 + |\boldsymbol{\xi}|^2)^{-\alpha}$:
\begin{align*}
\|\mathcal{R}_{m \to n}^\alpha f\|_{H^{s+\alpha-\frac{m-n}{2}}}^2 
&\le C \int_{\mathbb{R}^n} (1+|\boldsymbol{\xi}|^2)^{s+\alpha-\frac{m-n}{2}} (1+|\boldsymbol{\xi}|^2)^{-\alpha} (1+|\boldsymbol{\xi}|^2)^{-s+\frac{m-n}{2}} \notag \\
&\quad \times \left[\int_{\mathbb{R}^{m-n}} (1+|\mathbf{P}^T\boldsymbol{\xi}+\boldsymbol{\eta}|^2)^s |\widehat{f}|^2 d\boldsymbol{\eta}\right] d\boldsymbol{\xi} \\
&= C \int_{\mathbb{R}^m} (1 + |\mathbf{k}|^2)^s |\widehat{f}(\mathbf{k})|^2 d\mathbf{k} = C \|f\|_{H^s(\mathbb{R}^m)}^2,
\end{align*}
which establishes the sharp boundedness.
\end{proof}

\begin{corollary}[Critical Isomorphic Trace Parameter]
When the fractional order is chosen as $\alpha^* = \frac{m-n}{2}$, the Fractional Radon--Beta operator maps isomorphically between identical Sobolev regularity indices:
\begin{equation}
\mathcal{R}_{m \to n}^{\frac{m-n}{2}} : H^s(\mathbb{R}^m) \longrightarrow H^s(\mathbb{R}^n),
\end{equation}
achieving complete dimensional trace compensation without derivative loss.
\end{corollary}

\begin{theorem}[Dual Simplicial Extension Operator]
For dimensional elevation ($m < n$), the dual adjoint operator $\mathcal{E}_{m \to n}^\alpha \coloneqq (\mathcal{R}_{n \to m}^\alpha)^*$ is bounded:
\begin{equation}
\mathcal{E}_{m \to n}^\alpha : H^s(\mathbb{R}^m) \longrightarrow H^{s + \alpha + \frac{n-m}{2}}(\mathbb{R}^n).
\end{equation}
\end{theorem}

% ==============================================================================
\section{Coupled Multiscale 3D--2D--1D Interface Dynamical Systems}
% ==============================================================================

Consider a multiscale coupled physical system consisting of a 3D bulk domain $\Omega \subset \mathbb{R}^3$, an active 2D interface membrane $\Gamma \subset \mathbb{R}^2$, and a 1D conductive filament $\Sigma \subset \mathbb{R}^1$. Let $u(\mathbf{x}, t)$, $v(\mathbf{y}, t)$, and $w(z, t)$ denote the respective state fields.

\begin{definition}[Coupled Inter-Dimensional Transport System]
We formulate the transmission-coupled dynamical system:
\begin{equation}
\label{eq:coupled_system}
\begin{cases}
\partial_t u(\mathbf{x}, t) = \Delta u(\mathbf{x}, t) - \mathcal{E}_{2 \to 3}^\alpha\big(\lambda (\mathcal{R}_{3 \to 2}^\alpha u - v)\big)(\mathbf{x}, t), & \mathbf{x} \in \Omega \subset \mathbb{R}^3, \\
\partial_t v(\mathbf{y}, t) = \Delta_{\Gamma} v(\mathbf{y}, t) + \lambda \big(\mathcal{R}_{3 \to 2}^\alpha u - v\big)(\mathbf{y}, t) - \mathcal{E}_{1 \to 2}^\beta\big(\kappa (\mathcal{R}_{2 \to 1}^\beta v - w)\big)(\mathbf{y}, t), & \mathbf{y} \in \Gamma \subset \mathbb{R}^2, \\
\partial_t w(z, t) = \partial_z^2 w(z, t) + \kappa \big(\mathcal{R}_{2 \to 1}^\beta v - w\big)(z, t), & z \in \Sigma \subset \mathbb{R}^1,
\end{cases}
\end{equation}
where $\lambda, \kappa > 0$ are transmission coefficients, and $\alpha = 1/2, \beta = 1/2$ are the critical trace-compensating fractional orders.
\end{definition}

\begin{theorem}[Conservation of Total Mass and Energy Dissipation]
For any smooth solution of \eqref{eq:coupled_system}:
\begin{enumerate}
    \item \textbf{Conservation of Total Mass:} The total joint mass is strictly invariant for all $t \ge 0$:
    \begin{equation}
    \mathcal{M}_{\text{total}}(t) \coloneqq \int_\Omega u(\mathbf{x}, t)d\mathbf{x} + \int_\Gamma v(\mathbf{y}, t)d\mathbf{y} + \int_\Sigma w(z, t)dz = \mathcal{M}_{\text{total}}(0).
    \end{equation}
    \item \textbf{Global Energy Dissipation:} The joint energy $\mathcal{E}(t) \coloneqq \frac{1}{2}\|u\|_{L^2}^2 + \frac{1}{2}\|v\|_{L^2}^2 + \frac{1}{2}\|w\|_{L^2}^2$ dissipates monotonically:
    \begin{equation}
    \frac{d\mathcal{E}}{dt} = -\|\nabla u\|_{L^2(\Omega)}^2 - \|\nabla_\Gamma v\|_{L^2(\Gamma)}^2 - \|\partial_z w\|_{L^2(\Sigma)}^2 - \lambda \|\mathcal{R}_{3 \to 2}^\alpha u - v\|_{L^2(\Gamma)}^2 - \kappa \|\mathcal{R}_{2 \to 1}^\beta v - w\|_{L^2(\Sigma)}^2 \le 0.
    \end{equation}
\end{enumerate}
\end{theorem}

\begin{proof}
Integrating each equation over its domain, the transmission fluxes between bulk and interface satisfy:
\begin{equation}
\int_\Omega \mathcal{E}_{2 \to 3}^\alpha \big(\lambda (\mathcal{R}_{3 \to 2}^\alpha u - v)\big) d\mathbf{x} = \int_\Gamma \lambda (\mathcal{R}_{3 \to 2}^\alpha u - v) \mathcal{R}_{3 \to 2}^\alpha(1) d\mathbf{y} = \int_\Gamma \lambda (\mathcal{R}_{3 \to 2}^\alpha u - v) d\mathbf{y},
\end{equation}
since $\mathcal{R}_{3 \to 2}^\alpha(1) \equiv 1$ by Beta-integral normalization. Hence, the mass loss from $\Omega$ identically matches the gain on $\Gamma$. Similarly, the flux between $\Gamma$ and $\Sigma$ cancels, proving $\frac{d\mathcal{M}_{\text{total}}}{dt} = 0$.

For energy dissipation, multiplying by the state fields and applying the adjoint duality $\langle u, \mathcal{E}_{2 \to 3}^\alpha \phi \rangle_{L^2(\Omega)} = \langle \mathcal{R}_{3 \to 2}^\alpha u, \phi \rangle_{L^2(\Gamma)}$ with $\phi = \lambda(\mathcal{R}_{3 \to 2}^\alpha u - v)$:
\begin{align*}
-\langle u, \mathcal{E}_{2 \to 3}^\alpha \phi \rangle_{L^2(\Omega)} + \langle v, \phi \rangle_{L^2(\Gamma)} &= -\langle \mathcal{R}_{3 \to 2}^\alpha u, \phi \rangle_{L^2(\Gamma)} + \langle v, \phi \rangle_{L^2(\Gamma)} \\
&= -\langle \mathcal{R}_{3 \to 2}^\alpha u - v, \lambda(\mathcal{R}_{3 \to 2}^\alpha u - v) \rangle_{L^2(\Gamma)} = -\lambda \|\mathcal{R}_{3 \to 2}^\alpha u - v\|_{L^2(\Gamma)}^2 \le 0.
\end{align*}
Combining with the negative-definite Dirichlet integrals yields monotonic dissipation $\frac{d\mathcal{E}}{dt} \le 0$.
\end{proof}

% ==============================================================================
\section{Tomographic Reconstruction and Gibbs Artifact Elimination}
% ==============================================================================

Let $\operatorname{Gr}(n, m)$ denote the Grassmannian manifold of $n$-dimensional subspaces in $\mathbb{R}^m$. For projection matrix $\mathbf{P}_\theta$, let $g_\theta(\mathbf{z}) = \mathcal{R}_{m \to n, \theta}^\alpha f(\mathbf{z})$ be the observed simplicial projection.

\begin{theorem}[Exact Dual Inversion Formula and Gibbs Suppression]
For any test function $f \in \mathcal{S}(\mathbb{R}^m)$, $f$ is uniquely and stably reconstructed from $\{g_\theta\}_{\theta \in \operatorname{Gr}(n, m)}$ via:
\begin{equation}
f(\mathbf{x}) = \frac{1}{(2\pi)^m}\int_{\operatorname{Gr}(n, m)} d\mu(\theta) \int_{\mathbb{R}^n} e^{i \boldsymbol{\xi}\cdot \mathbf{P}_\theta \mathbf{x}} |\boldsymbol{\xi}|^{m-n} \frac{\widehat{g}_\theta(\boldsymbol{\xi})}{\widehat{\mathcal{K}}_\alpha(\mathbf{P}_\theta^T \boldsymbol{\xi})} \, d\boldsymbol{\xi}.
\end{equation}
Furthermore, because $\widehat{\mathcal{K}}_\alpha(\mathbf{k})$ decays smoothly as $\mathcal{O}(|\mathbf{k}|^{-\alpha})$ without discontinuity, reconstruction of discontinuous edge profiles avoids Gibbs ringing phenomena, providing unconditional $L^2$-error bounds.
\end{theorem}

% ==============================================================================
\section{Simplicial Geometry-Preserving Dimensionality Reduction on Hypergraphs}
% ==============================================================================

In high-dimensional data analysis, multi-way relational structures are encoded as abstract simplicial complexes $K = \{\sigma_1, \dots, \sigma_N\}$ with barycentric coordinates $\mathbf{x} = \sum_{j=1}^m c_j \mathbf{v}_j$, where $\sum c_j = 1, c_j \ge 0$.

\begin{theorem}[Barycentric Ratio Preservation]
Let $\sigma \in K$ be an $(m-1)$-simplex embedded in $\mathbb{R}^m$. The operator $\mathcal{R}_{m \to n}^\alpha$ projects $\sigma$ into an $(n-1)$-simplex $\tau \subset \mathbb{R}^n$ such that for any two interior points $\mathbf{x}_1, \mathbf{x}_2 \in \sigma$:
\begin{equation}
\frac{\operatorname{dist}_{\Delta}(\mathcal{R}_{m \to n}^\alpha \mathbf{x}_1, \mathcal{R}_{m \to n}^\alpha \mathbf{x}_2)}{\operatorname{dist}_{\Delta}(\mathbf{x}_1, \mathbf{x}_2)} = 1 + \mathcal{O}(\alpha^{-1}),
\end{equation}
where $\operatorname{dist}_{\Delta}$ denotes the barycentric Wasserstein metric.
\end{theorem}

% ==============================================================================
\section{Grassmannian Matrix-Valued Beta-Kernels}
% ==============================================================================

Let $\operatorname{Sym}_m^+(\mathbb{R})$ denote the open convex cone of $m \times m$ real symmetric positive-definite matrices.

\begin{definition}[Siegel--Wishart Matrix Beta Operator]
For matrix parameters $\mathbf{A}, \mathbf{B} \in \operatorname{Sym}_m^+(\mathbb{R})$ with $\mathbf{A}, \mathbf{B} > \frac{m-1}{2}\mathbf{I}_m$, the matrix Beta operator acting on $F: \operatorname{Sym}_m^+(\mathbb{R}) \to \mathbb{R}$ is defined by
\begin{equation}
\mathcal{G}_{\mathbf{A}, \mathbf{B}} F(\mathbf{X}) \coloneqq \frac{1}{\mathrm{B}_m(\mathbf{A}, \mathbf{B})} \int_{\mathbf{0} < \mathbf{Y} < \mathbf{I}_m} (\det \mathbf{Y})^{\mathbf{A} - \frac{m+1}{2}} (\det(\mathbf{I}_m - \mathbf{Y}))^{\mathbf{B} - \frac{m+1}{2}} F\left(\mathbf{X}^{1/2} \mathbf{Y} \mathbf{X}^{1/2}\right) d\mathbf{Y},
\end{equation}
where the multivariate Siegel Beta normalization is given by
\begin{equation}
\mathrm{B}_m(\mathbf{A}, \mathbf{B}) \coloneqq \frac{\Gamma_m(\mathbf{A})\Gamma_m(\mathbf{B})}{\Gamma_m(\mathbf{A}+\mathbf{B})}, \quad \Gamma_m(\mathbf{A}) \coloneqq \pi^{m(m-1)/4}\prod_{j=1}^m \Gamma\left(\mathbf{A} - \frac{j-1}{2}\right).
\end{equation}
\end{definition}

\begin{theorem}[Invariance and Spherical Harmonic Eigenfunctions]
The operator $\mathcal{G}_{\mathbf{A}, \mathbf{B}}$ is invariant under the congruence action of the orthogonal group $\mathrm{O}(m)$ ($\mathbf{X} \mapsto \mathbf{U}\mathbf{X}\mathbf{U}^T$). Furthermore, for any zonal spherical polynomial $Z_\lambda(\mathbf{X})$ associated with partition $\lambda$:
\begin{equation}
\mathcal{G}_{\mathbf{A}, \mathbf{B}} Z_\lambda(\mathbf{X}) = \frac{[\mathbf{A}]_\lambda}{[\mathbf{A}+\mathbf{B}]_\lambda} Z_\lambda(\mathbf{X}),
\end{equation}
where $[\mathbf{A}]_\lambda = \prod_{j=1}^m \frac{\Gamma(a_j + \lambda_j)}{\Gamma(a_j)}$ denotes the multivariate generalized Pochhammer symbol.
\end{theorem}

% ==============================================================================
\section{Conclusion}
% ==============================================================================

We have established a comprehensive functional-analytic theory of dimension-changing non-local operators $\mathcal{T}_{m \to n}^\alpha$. We proved the sharp Sobolev trace mapping theorem, formulated conservative coupled 3D--2D--1D interface dynamical systems, demonstrated Gibbs-free tomographic reconstruction, established barycentric preservation in hypergraph data science, and generalized the theory to matrix Grassmannians. Future work will investigate numerical finite-element schemes for inter-dimensional multiscale PDEs.

\begin{thebibliography}{99}
\bibitem{silvafilho2026simplex} Reinaldo Maia Silva-Filho, \emph{Analytic Continuation of Pascal's Simplex: Continuous Multinomial Integrals, Polytope Boundary Recurrences, and Simplicial Fractional Calculus}, Paper I of the Series, 2026.
\bibitem{silvafilho2026waves} Reinaldo Maia Silva-Filho, \emph{Nonlinear Simplicial Waves and Anomalous Porous Transport Induced by Beta-Kernel Fractional Laplacians}, Paper II of the Series, 2026.
\bibitem{helgason2011} S. Helgason, \emph{Integral Geometry and Radon Transforms}, Springer, New York, 2011.
\bibitem{samko1993} S. G. Samko, A. A. Kilbas, and O. I. Marichev, \emph{Fractional Integrals and Derivatives: Theory and Applications}, Gordon and Breach, Yverdon, 1993.
\bibitem{adams2003} R. A. Adams and J. J. F. Fournier, \emph{Sobolev Spaces}, Academic Press, New York, 2003.
\bibitem{muirhead2009} R. J. Muirhead, \emph{Aspects of Multivariate Statistical Theory}, John Wiley \& Sons, Hoboken, 2009.
\bibitem{evans2010} L. C. Evans, \emph{Partial Differential Equations}, American Mathematical Society, Providence, 2010.
\end{thebibliography}

\end{document}
